% !TeX spellcheck = pl

\documentclass[a4paper]{report}

\usepackage{baskervald}  % Default font
\usepackage{cabin}  % Sans serif font
\usepackage[T1]{fontenc}
\usepackage[utf8]{inputenc}
\usepackage{setspace} \onehalfspacing
\usepackage{graphicx}
\usepackage{textcomp}
\usepackage{hyperref}

% This macro produces a "Snap!" logo.
%
% Note that in Polish (and other languages), nouns are inflected. The form "Snap!" with suffix looks plain stupid, so the macro takes a suffix as argument. If the suffix is non-empty, no exclamation mark is generated, and the suffix is used instead.
\newcommand{\Snap}[1]{%
	\textsf{%
		Snap%
		\ifx\\#1\\\textit{!\@}%
		\else #1%
		\fi%
	}%
}

\newcommand{\code}[1]{\textsf{#1}}
\newcommand{\defaultGraphicsScale}{0.6}

\renewcommand{\thechapter}{\Roman{chapter}}
\renewcommand{\thesection}{\Alph{section}}
\renewcommand{\thesubsection}{}

\newcommand{\inlinepic}[1]{%
	\raisebox{-4pt}{%
		\includegraphics[scale=\defaultGraphicsScale]{#1}%
	}%
}

\newcommand{\bigpic}[1]{
	\begin{center}
	\includegraphics[scale=\defaultGraphicsScale]{#1}
	\end{center}
}

\usepackage[polish]{babel}
\usepackage{polski}
\frenchspacing
\usepackage{indentfirst}

\begin{document}

\title{Snap! --- Podręcznik użytkownika}
\author{Brian Harvey\texorpdfstring{ \and}{,} Jens M\"onig}
\date{}

\manualtitlepage[Tłumaczenie na język polski:\\Bartosz Leper]{Snap!\\Podręcznik użytkownika}

\tableofcontents

\chapter*{}
\section*{Podziękowania}

Mieliśmy ogromne szczęście do mentorów. Jens zdobył dużo doświadczenia pracując wśród pionierów Smalltalka: Alana Kaya, Dana Ingallsa i~reszty ekipy, która wynalazła komputery osobiste i~programowanie obiektowe w~najlepszych dniach firmy Xerox PARC. Pracował z~Johnem Maloneyem z~zespołu Scratcha w~MIT\footnote{Massachusetts Institute of Technology, amerykańska uczelnia techniczna --- przyp. tłum.}, autorem platformy graficznej Morphic, wciąż stanowiącej fundament \Snap{a}. Znakomity projekt języka Scratch, autorstwa Lifelong Kindergarten Group z~MIT Media Lab, odgrywa w~\Snap{ie} kluczową rolę.

\textbf{\emph{Nasza poprzednia wersja, BYOB, była bezpośrednią modyfikacją kodu źródłowego Scratcha. \Snap{} został napisany od zera, lecz struktura jego kodu oraz interfejs użytkownika pozostają mocno zakorzenione w~Scratchu. Z~kolei zespół Scratcha, który mógłby widzieć w~nas rywali, przyjął nas ciepło i~okazał nam całkowite wsparcie.}}

Brian zdobywał szlify w~MIT oraz Stanford Artificial Intelligence Labs\footnote{Laboratorium sztucznej inteligencji na Uniwersytecie Stanforda --- przyp. tłum.}, gdzie uczył się pod okiem Johna McCarthy'ego, twórcy Lispa, oraz Geralda~J. Suss\-mana i~Guya Steele'a, twórców języka Scheme. Zdobywał również wiedzę od wielu innych wybitnych informatyków, w~tym autorów najlepszej książki z zakresu informatyki --- \emph{Struktury i~interpretacji programów komputerowych}: Hala Abelsona, Geralda~J. Suss\-mana i~Julie Suss\-man.

\textbf{\emph{Za starych dobrych czasów mawialiśmy w~MIT Logo Lab: ,,Język Logo to Lisp w przebraniu BASIC-a''. Dziś, ze swoimi pierwszoklasowymi procedurami, zasięgami leksykalnymi~i pierwszoklasowymi kontynuacjami, \Snap{} jest jak Scheme w~przebraniu Scratcha.}}

Szczęśliwym zrządzeniem losu, poprzez forum Scratch Advanced Topics, poznaliśmy wspaniałą grupę błyskotliwych uczniów gimnazjów~(!\@) i liceów. Kilku z nich wniosło swój wkład w~kod \Snap{a}: Kartik Chandra, Nathan Dinsmore, Connor Hudson i~Ian Reynolds. Ponadto wielu zgłosiło pomysły i~raporty błędów podczas testowania wersji alfa. Wśród studentów Uniwersytetu Kalifornijskiego w~Berkeley, którzy przyczynili się do rozwoju kodu, znajdują się Michael Ball, Achal Dave, Kyle Hotchkiss, Ivan Motyashov i~Yuan Yuan. Wymienianie wszystkich tłumaczy zajęłoby zbyt wiele miejsca, ale można ich odnaleźć w~okienku ,,O \Snap{}\ldots'' dostępnym w~programie. Niniejsze dzieło powstało częściowo w~ramach grantu nr~1143566 udzielonego przez National Science Foundation, a częściowo przy wsparciu firmy MioSoft.

\clearpage

\begin{center}
\bf \Huge \Snap{} \\
Podręcznik użytkownika \\
\huge Wersja 4.0 \vspace{40pt}
\end{center}

\Snap{} to rozszerzona reimplementacja języka Scratch (\url{http://scratch.mit.edu}), która pozwala na tworzenie własnych bloków (ang.\ \textit{Build Your Own Blocks}; stąd dawna nazwa \Snap{a} --- BYOB). Opisywany tu język obsługuje pierwszoklasowe listy, procedury i~kontynuacje. Te dodatkowe możliwości sprawiają, że nadaje się on do przeprowadzenia poważnego wstępu do informatyki dla uczniów liceów i szkół wyższych. Aby uruchomić środowisko \Snap{}, wystarczy otworzyć przeglądarkę internetową i~wpisać adres \url{http://snap.berkeley.edu/run}, aby zacząć pracę z~minimalnym zestawem bloków. Można też użyć adresu \url{http://snap.berkeley.edu/init}, aby załadować niewielki zestaw dodatkowych bloków. Wiąże się to z~nieco wolniejszym ładowaniem, ale jest zalecane dla wygody użytkowników (w~dalszej części podręcznika będziemy zakładali korzystanie z~tej właśnie metody).

\clearpage

\chapter{Bloki, skrypty i~duszki}

W~tym rozdziale poznamy kilka cech języka \Snap{} odziedziczonych po Scratchu; doświadczeni użytkownicy Scratcha mogą przejść od razu do sekcji~\ref{sec:zagnieżdżanie-duszków}.

\Snap{} jest językiem programowania --- notacją, przy pomocy której możemy powiedzieć komputerowi, co ma zrobić. Jednak w~odróżnieniu od większości innych, \Snap{} jest językiem wizualnym; programując w~nim, zamiast posługiwać się klawiaturą, używamy metody ,,przeciągnij i~upuść'', dobrze znanej użytkownikom komputerów.

Uruchom teraz środowisko \Snap{}. Powinieneś zobaczyć ekran podzielony na kilka obszarów:\footnote{\onehalfspacing Pierwsze uruchomienie \Snap{a} prawdopodobnie spowoduje wyświetlenie angielskiej wersji programu; aby przełączyć się na język polski, należy kliknąć menu ,,Ustawienia'' \inlinepic{../common/btn-settings} na pasku narzędzi, a~następnie użyć polecenia ,,Language\ldots'' (,,Język\ldots'') --- przyp. tłum.}\nopagebreak

\begin{center}
\def\svgwidth{\textwidth}
\input{obszary-okna.pdf_tex}
\end{center}

(Proporcje tych stref mogą się różnić, w~zależności od rozmiaru i~kształtu okna przeglądarki).

Program w~języku \Snap{} składa się z~jednego lub więcej \emph{skryptów}, te zaś z~kolei --- z~\emph{bloków}. Oto przykładowy skrypt:\nopagebreak

\label{fig:typowy-skrypt}
\bigpic{typowy-skrypt}

Na powyższy skrypt składa się pięć bloków w~trzech różnych kolorach, odpowiadających trzem z~ośmiu \emph{palet} z~blokami. Obszar palet, znajdujący się po lewej stronie okna, pokazuje jedną paletę na raz. Do zmiany widocznej palety służy osiem przycisków znajdujących się tuż nad tym obszarem. Bloki ciemnożółte, widoczne w~naszym skrypcie, pochodzą z~palety ,,Kontrola''; zielone z~palety ,,Pisak'', a~niebieskie --- z~palety ,,Ruch''. Aby złożyć taki skrypt, należy poprzeciągać odpowiednie bloki z~palet do \emph{obszaru skryptów}, umiejscowionego na środku okna. Kiedy układamy jeden blok pod drugim w~taki sposób, aby wcięcie dolnego bloku znalazło się w~pobliżu wypustki tego powyżej, bloki łączą się ze sobą (ang. \textit{snap together}; stąd nazwa języka \Snap{}):\nopagebreak

\bigpic{laczenie-blokow}

Pozioma biała linia sygnalizuje, że jeśli puścimy zielony blok, połączy się on z~wypustką ciemnożółtego.

\subsection{Bloki-czapki i~bloki komend}

Na górze skryptu znajduje się \emph{blok-czapka}, który określa, kiedy skrypt ma zostać wykonany. Nazwy bloków-czapek zazwyczaj zaczynają się słowem ,,\code{kiedy}''; nasz przykładowy skrypt powinien zostać uruchomiony w~momencie kliknięcia zielonej flagi, znajdującej się w pobliżu prawej strony paska narzędzi \Snap{a}. (Pasek ten jest częścią okna programu \Snap{}; nie chodzi tutaj o pasek menu przeglądarki lub systemu operacyjnego). Skrypt nie musi posiadać czapki, jednak w~takim przypadku zostanie wykonany tylko wtedy, gdy użytkownik sam go kliknie. Skrypt nie może mieć więcej niż jednej czapki; jej charakterystyczny kształt służy łatwiejszemu zapamiętaniu tej szczególnej własności.

Pozostałe bloki w naszym skrypcie to \emph{bloki komend}. Każdy z~nich oznacza jakąś akcję, którą \Snap{} potrafi wykonać. Na przykład blok \inlinepic{przesun-o-10-krokow} nakazuje duszkowi\footnote{W grafice komputerowej słowem ,,duszek'' (ang. \textit{sprite}) nazywa się ruchomy obiekt na ekranie --- przyp. tłum.}, czyli strzałce na \emph{scenie} po prawej stronie okna, aby przesunął się o~dziesięć kroków do przodu w~kierunku, w~którym jest zwrócony. Każdy krok to niewielka odległość na ekranie. Wkrótce przekonamy się, że na scenie może być więcej duszków, a~każdy z nich może mieć własne skrypty. Ponadto duszki nie muszą wyglądać jak strzałki; ich kostiumy mogą być dowolnymi obrazkami. Kształt bloku \code{przesuń} ma za zadanie przypominać klocek, skrypt zaś jest jak wieża z klocków. Słowa ,,blok'' będziemy używać dla oznaczenia zarówno graficznego symbolu na ekranie, jak i~procedury (akcji) jaką ten blok wykonuje.

Liczbę $10$ w powyższym bloku \code{przesuń} nazywamy jego \emph{parametrem}. Kliknąwszy na białym, owalnym polu, możemy wpisać w~jej miejsce dowolną inną. W przykładowym skrypcie ze strony \pageref{fig:typowy-skrypt} parametrem jest liczba $100$. Jak się później okaże, pola parametrów mogą mieć kształty inne od owalnych; oznacza to wtedy, że akceptują one wartości inne niż liczby. Zobaczymy również, że zamiast wpisywać konkretne wartości w~pola, możemy nakazać komputerowi je obliczać. Ponadto blok może mieć więcej niż jeden parametr. Na przykład blok \code{leć}, znajdujący się mniej więcej w~połowie palety ,,Ruch'', przyjmuje trzy parametry.

Większość bloków komend ma kształt klocków, lecz niektóre, jak \code{powtórz} z~tego samego przykładu, wyglądają jak \emph{klamry}. Większość bloków klamrowych można znaleźć na palecie ,,Kontrola'. Wnętrze klamry jest szczególnym rodzajem pola parametru, który przyjmuje \emph{skrypt} jako parametr. W~przykładowym skrypcie blok \code{powtórz} ma dwa parametry: liczbę $4$ oraz skrypt\nopagebreak

\bigpic{typowy-skrypt-wnetrze}

\section{Duszki i~współbieżność}

Tuż pod sceną znajduje się przycisk ,,nowy duszek'' \inlinepic{../common/btn-new-sprite}. Kliknięcie go spowoduje dodanie nowego duszka do sceny. Pojawi się on w~losowym miejscu na scenie, skierowany w~losową stronę i~zabarwiony na losowy kolor.

Każdy duszek ma swoje własne skrypty. Aby wyświetlić w~obszarze skryptów te należące do konkretnego duszka, należy kliknąć na jego obrazku w~\emph{zagrodzie duszków}, znajdującej się w~prawym dolnym rogu okna. Spróbuj umieścić następujące skrypty w~obszarze skryptów --- po jednym dla każdego duszka:\nopagebreak

\begin{figure}[H]
\begin{minipage}{0.5\textwidth}
\centering
\includegraphics[scale=\defaultGraphicsScale]{duszki-i-wspolbieznosc-1}
\end{minipage}%
\begin{minipage}{0.5\textwidth}
\centering
\includegraphics[scale=\defaultGraphicsScale]{duszki-i-wspolbieznosc-2}
\end{minipage}
\end{figure}

Kiedy klikniemy zieloną flagę \inlinepic{../common/btn-start}, powinniśmy zobaczyć jak jeden duszek się obraca, podczas gdy drugi porusza się w~tę i~z~powrotem. Ten eksperyment pokazuje, jak różne skrypty mogą być wykonywane jednocześnie (\emph{współbieżnie}). Obracanie się dookoła i~ruch po linii prostej zachodzą jednocześnie. Współbieżność zachodzi również w~przypadku wielu skryptów należących do tego samego duszka. Spróbujmy tego przykładu:\nopagebreak

\begin{figure}[H]
\begin{minipage}{0.5\textwidth}
\centering
\includegraphics[scale=\defaultGraphicsScale]{duszki-i-wspolbieznosc-3}
\end{minipage}%
\begin{minipage}{0.5\textwidth}
\centering
\includegraphics[scale=\defaultGraphicsScale]{duszki-i-wspolbieznosc-4}
\end{minipage}
\end{figure}

Po naciśnięciu spacji duszek powinien zacząć bez końca chodzić w kółko, ponieważ bloki \code{przesuń} i \code{obróć} są wykonywane współbieżnie. (Aby przerwać program, kliknij czerwony czerwony znak ,,stop'' \inlinepic{../common/btn-stop} na prawym brzegu paska narzędzi).

\subsection{Kostiumy i~dźwięki}

Aby zmienić wygląd duszka, należy zaimportować dla niego nowy \emph{kostium}. Są na to trzy sposoby. Najpierw trzeba wybrać duszka z~zagrody. Następnie, w~pierwszej metodzie, klikamy na ikonie pliku \inlinepic{../common/btn-file} na pasku narzędzi, a~następnie wybieramy polecenie ,,\code{Kostiumy\ldots}''. Ukaże się lista kostiumów z~publicznej biblioteki multimediów, spośród których możemy dokonać wyboru. Drugą metodą jest wybór pliku ze swojego własnego kompuera. Należy w~tym celu kliknąć ikonę pliku, a~następnie polecenie ,,\code{Importuj\ldots}''. Można wtedy wybrać plik obrazu w~dowolnym formacie (PNG, JPEG itd.) obsługiwanym przez przeglądarkę. Trzeci sposób jest szybszy jeśli plik, którego chcemy użyć, jest widoczny na pulpicie: po prostu przeciągnij go do okna \Snap{a}. W~każdym z~tych przypadków obszar skryptów zacznie wyglądać mniej więcej tak:\nopagebreak

\bigpic{obszar-skryptow-z-dodatkowym-kostiumem}

Tuż nad tą częścią okna znajdują się trzy zakładki: ,,Skrypty'', ,,Kostiumy'' i~,,Dźwię\-ki''. W~tym momencie aktywna jest karta ,,Kostiumy''. Widzimy na niej \emph{garderobę} duszka i~możemy z~jej poziomu wybrać dla niego kostium --- domyślny kostium żółwia\footnote{Z powodów historycznych, słowem ,,żółw'' nazywamy ruchomy obiekt, który porusza się wykonując program i~rysuje, zostawiając za sobą ślad} lub wybrany wcześniej kostium Alonza. (Alonzo, maskotka \Snap{a}, został nazwany na cześć Alonza Churcha, matematyka, który jako pierwszy wpadł na pomysł, aby procedury traktować na równi z~danymi, co jest najistotniejszą różnicą między \Snap{em} a~Scratchem). Możemy przypisać duszkowi tyle kostiumów ile chcemy, a~potem wybierać, który z~nich założy, albo poprzez kliknięcie w~obrębie garderoby, albo używając w~skrypcie bloku \inlinepic{zmien-kostium-na-zolwia} lub \inlinepic{nastepny-kostium}. (Każdy kostium ma zarówno numer jak i~nazwę. Blok \code{następny kostium} wybiera następny w~kolejności kostium; po ostatnim wybiera z~powrotem kostium numer~1. Żółw, czyli kostium numer~0, jest przez blok \code{następny kostium} ignorowany). Kostium ,,Żółw'' jest jedynym, który zmienia kolor zgodnie z~kolorem pisaka.

Oprócz kostiumów, duszki mogą posiadać \emph{dźwięki}; dźwiękowy odpowiednik garderoby duszka nazywamy jego \emph{szafą grającą}. Można importować pliki dźwiękowe w~dowolnym formacie obsługiwanym przez przeglądarkę. Do odtwarzania dźwięków służą dwa rodzaje bloków. Jeśli skrypt ma się dalej wykonywać podczas odtwarzania, używamy bloku \inlinepic{zagraj-dzwiek-ratunku}. Za to aby poczekać, aż dźwięk się zakończy, zanim skrypt będzie kontynuowany, należy wykorzystać blok \inlinepic{zagraj-dzwiek-ratunku-i-czekaj}.

\subsection{Nadawanie i~odbieranie komunikatów}

Widzieliśmy wcześniej przykład dwóch duszków poruszających się jednocześnie. Jednak w~bardziej interesującym programie duszki na scenie będą wchodzić w~interakcje, abyśmy mogli opowiedzieć przy ich pomocy jakąś historię, zagrać w~grę itd. Czasami jeden duszek będzie musiał nakazać innemu wykonanie jakiegoś skryptu. Oto prosty przykład:\nopagebreak

\begin{figure}[H]
\begin{minipage}{0.5\textwidth}
\centering
\includegraphics[scale=0.4]{../common/boy1-walking}
\end{minipage}%
\begin{minipage}{0.5\textwidth}
\centering
\reflectbox{\includegraphics[scale=0.3]{../common/dog2-c}}
\end{minipage}
\vskip 3ex
\begin{minipage}[t]{0.5\textwidth}
\centering
\vspace{0pt} % REALLY align to top
\includegraphics[scale=\defaultGraphicsScale]{nadawanie-i-odbieranie-komunikatow-1}
\end{minipage}%
\begin{minipage}[t]{0.5\textwidth}
\centering
\vspace{0pt} % REALLY align to top
\includegraphics[scale=\defaultGraphicsScale]{nadawanie-i-odbieranie-komunikatow-2}
\end{minipage}
\end{figure}

Słowo ,,szczekaj'' występujące w~bloku \inlinepic{nadaj-szczekaj-do-wszystkich-i-czekaj} to pierwszy lepszy wyraz, który przyszedł mi do głowy. Jedną z~opcji, które ukazują się po kliknięciu strzałki w~dół obok tego pola parametru (i~jedyną początkowo dostępną), jest ,,\code{nowy}''. Po jej wybraniu \Snap{} pyta o~nazwę komunikatu. Kiedy wspomniany blok zostanie wykonany, wybrany komunikat zostaje wysłany do \emph{każdego} duszka --- stąd też określenie ,,nadaj do wszystkich''. Jednak w~naszym przykładzie tylko jeden duszek ma skrypt, który jest wywoływany w~momencie nadania tego komunikatu --- jest nim pies. Ponieważ skrypt chłopca wykorzystuje blok \code{nadaj do wszystkich i~czekaj} zamiast \code{nadaj do wszystkich}, chłopiec nie przechodzi do następującego po nim bloku \code{powiedz}, dopóki skrypt psa się nie skończy. Z~tej przyczyny dwa duszki mówią na zmianę, a~nie jednocześnie.

Warto przy okazji zwrócić uwagę na to, że pierwsze pole parametru na bloku \code{powiedz} nie jest owalne, lecz prostokątne. Oznacza to, że parametr ten może być dowolnym łańcuchem znaków (tekstem), nie tylko liczbą. W~polach parametrów typu tekstowego spacje ukazują się jako brązowe kropki, abyśmy mogli policzyć liczbę odstępów między wyrazami. Co ważniejsze, możemy dzięki temu odróżnić pusty łańcuch od złożonego z~samych spacji. Brązowe kropki \emph{nie będą} widoczne na scenie, kiedy blok zostanie wykonany.

Scena ma swój własny obszar skryptów. Możemy wyświetlić jej szczegóły klikając ikonę ,,Scena'' po lewej stronie zagrody duszków. W~przeciwieństwie do duszków scena się nie porusza. Zamiast kostiumów ma \emph{tła} --- obrazki wypełniające cały obszar sceny. Duszki rysowane są na aktualnym tle. W~skomplikowanych projektach często wygodnie jest użyć skryptu sceny do koordynacji działań poszczególnych części programu.

\section{Zagnieżdżanie duszków: kotwice i~części}
\label{sec:zagnieżdżanie-duszków}

Czasem dobrze jest stworzyć swego rodzaju ,,nadduszka'', złożonego z~kawałków, które poruszają się razem, ale mogą być osobno względem siebie ustawiane. Klasycznym przykładem może być ciało człowieka złożone z~tułowia, kończyn i~głowy. \Snap{} pozwala nam uczynić jednego z~duszków \emph{kotwicą} złożonego obiektu, a~resztę --- jego \emph{częściami}. Aby zagnieździć w~ten sposób duszki, należy przeciągnąć z~zagrody ikonę duszka, który ma zostać \emph{częścią} złożonego obiektu na znajdującego się na scenie (nie w~zagrodzie!) duszka, który zostanie \emph{kotwicą}.

Zagnieżdżone duszki --- zarówno kotwice jak i części --- mają w zagrodzie specjalne oznaczenia:\nopagebreak

\bigpic{zagniezdzone-duszki-w-zagrodzie}

W~tym przypadku chcielibyśmy animować rękę Alonza. (Ręka została pokolorowana na zielono, aby uwypuklić zależność między dwoma duszkami, choć w~prawdziwym projekcie miałyby one raczej ten sam kolor). ,,Duszek1'', reprezentujący ciało Alonza, jest kotwicą; ,,Duszek2'' to ręka. Ikona duszka-kotwicy zawiera w~dolnej części do trzech miniatur doczepionych do niego duszków-części. Z~kolei na ikonie każdej z~części widać pomniejszony obrazek duszka-kotwicy w~lewym górnym rogu, w~prawym górnym zaś --- \emph{przełącznik synchronizacji obrotów}. Początkowo, jak widać na rysunku powyżej, jest on tak ustawiony, aby obrót kotwicy powodował zarówno orbitowanie części wokół niej, jak i~obrót części dookoła swojej własnej osi. Po kliknięciu przełącznik zmienia kształt z~okrągłej strzałki na prostą, co oznacza, że od tej pory obrót duszka-kotwicy będzie powodował jedynie zmianę pozycji przymocowanych do niego części, ale nie będą się one obracać wokół własnej osi. (Części mogą również obracać się niezależnie, przy pomocy bloków \code{obróć}). Każda zmiana pozycji lub rozmiaru kotwicy jest propagowana na wszystkie części.

\begin{figure}[H]
\centering
\includegraphics[scale=\defaultGraphicsScale]{komenda-machania-reka}%
\hspace{2em}%
\includegraphics[scale=0.4]{../common/alonzo-waving}
\end{figure}

\section{Bloki funkcji i~wyrażenia}

Jak dotąd używaliśmy dwóch rodzajów bloków: ,,czapek'' i~komend. Kolejnym rodzajem jest blok \emph{funkcji}, który ma owalny kształt: \inlinereporterpic{pozycja-x}. Nazywamy go ,,blokiem funkcji'', ponieważ --- podobnie jak funkcja w~matematyce --- kiedy zostaje wykonany, zamiast przeprowadzać jakąś czynność, zwraca wartość, która może zostać użyta jako parametr w~innym bloku. Jeśli przeciągniemy sam blok funkcji do obszaru skryptów i~klikniemy go, obok pokaże się dymek z~wartością zwróconą przez tę funkcję:\nopagebreak

\bigpic{pozycja-x-zwraca-liczbe}

Kiedy przeciągamy blok funkcji nad polem parametru należącym do innego bloku, wokół tego pola pojawia się biała otoczka, analogicznie do sytuacji, w~której łączymy bloki komend i~pojawia się biała linia. Oto przykładowy skrypt wykorzystujący funkcję:\nopagebreak

\begin{figure}[H]
\centering
\includegraphics[scale=\defaultGraphicsScale]{przykladowy-skrypt-wykorzystujacy-funkcje}%
\hspace{2em}%
\includegraphics{../common/turtle-says-its-position}
\end{figure}

Funkcja \code{pozycja X} nadaje tu wartość pierwszemu parametrowi bloku \code{powiedz}. Pozycja X~duszka to inaczej jego współrzędna pozioma. Określa ona, jak daleko w~lewo (jeśli jest liczbą ujemną) lub w~prawo (jeśli dodatnią) znajduje się duszek w~stosunku do środka sceny. Analogicznie, pozycja Y~to współrzędna pionowa, mierzona ilością kroków w~górę (wartości dodatnie) lub w~dół od środka (wartości ujemne).

Przy pomocy funkcji z palety ,,Wyrażenia'' możemy wykonywać obliczenia:\nopagebreak

\begin{figure}[H]
\centering
\includegraphics[scale=\defaultGraphicsScale]{uzycie-funkcji-do-obliczen}%
\hspace{2em}%
\includegraphics{../common/turtle-says-its-rounded-position}
\end{figure}

Blok \code{zaokrąglij} zaokrągla $35.3905\ldots$ do $35$, a~blok~\code{+} dodaje do tej liczby $100$. Nawiasem mówiąc, choć blok \code{zaokrąglij} znajduje się na palecie ,,Wyrażenia'', podobnie jak~\code{+}, to w~tym skrypcie ma on jaśniejszy kolor i~czarne litery. To dlatego, że \Snap{} używa na przemian ciemnych i~jasnych odcieni kolorów, kiedy zagnieżdżamy w~sobie bloki z~tej samej palety:\nopagebreak

\bigpic{kolorowanie-w-zebre}

Takie \emph{kolorowanie w~zebrę} poprawia czytelność programu. Blok funkcji wraz z~parametrami, a~być może również innymi blokami funkcji, na przykład \inlinepic{zaokraglij-pozycja-x-plus-100}, nazywamy \emph{wyrażeniem}.

\section{Predykaty i~obliczenia warunkowe}

Większość funkcji zwraca albo liczbę, jak \inlinereporterpic{plus}, lub łańcuch tekstowy, jak \inlinereporterpic{polacz-witaj-swiecie}. \emph{Predykat} to specjalny rodzaj funkcji, która zawsze zwraca jedną z dwojga wartości \code{prawdę} lub \code{fałsz}. Predykaty mają kształt sześciokątów:\nopagebreak

\bigpic{przycisk-myszy-nacisniety}

\begin{sloppypar}
Specjalny kształt jest oznaką, że predykaty nie mają z~reguły sensu w~tych polach parametrów, które oczekują liczby lub tekstu. Raczej nie napisalibyśmy \inlinepic{przesun-o-przycisk-myszy-nacisniety-krokow}, choć gdybyśmy się uparli, \Snap{} by nam na to pozwolił, co widać na załączonym obrazku. W~typowych sytuacjach predykaty umieszczamy w~specjalnych sześciokątnych polach parametrów takich jak to:\nopagebreak
\end{sloppypar}

\bigpic{jezeli-to}

Klamra \code{jeżeli --- to} wykonuje obejmowany przez nią fragment skryptu wtedy i~tylko wtedy, gdy wyrażenie w~jej sześciokątnym polu parametru jest prawdziwe, czyli zwraca wartość \code{prawda}.\nopagebreak

\bigpic{predykaty-i-obliczenia-warunkowe-1}

Poniższy blok jest bardzo użyteczny w~animacjach. Wykonuje on skrypt będący jego parametrem \emph{wielokrotnie}, dopóki predykat nie zostanie spełniony:\nopagebreak

\bigpic{predykaty-i-obliczenia-warunkowe-2}

Jeśli pracując nad projektem, będziemy chcieli tymczasowo pominąć niektóre komendy w~skrypcie, lecz nie będziemy chcieli zapomnieć, gdzie było ich miejsce, możemy użyć następującej sztuczki:\nopagebreak

\bigpic{predykaty-i-obliczenia-warunkowe-3}

Czasami potrzeba wykonać tę samą czynność bez względu na to, czy jakiś warunek zachodzi czy nie, za to z~różnymi parametrami dla obu tych przypadków. Można do tego użyć bloku \emph{funkcji} \code{if}:\footnote{\onehalfspacing Jeśli nie widzisz go w~palecie ,,Kontrola'', kliknij przycisk ,,Plik'' \inlinepic{../common/btn-file} na pasku narzędzi i~wybierz polecenie ,,Importuj narzędzia''.}\footnote{Niestety, podobnie jak pozostałe dodatkowe narzędzia i~biblioteki bloków, blok \code{if --- then --- else} posiada wyłącznie angielską nazwę, bez względu na nasze ustawienia języka. Oznacza ona ,,jeżeli --- to --- w~przeciwnym razie'' --- przyp. tłum.}\nopagebreak

\bigpic{predykaty-i-obliczenia-warunkowe-4}

\section{Zmienne}

Wypróbujmy następujący skrypt:\footnote{Blok \code{for}\footnotemark{} również znajduje się w~bibliotece narzędzi; użyj polecenia ,,\code{Importuj narzędzia}'' z~menu ,,Plik'', jeśli nie widzisz ich na palecie ,,Kontrola''.}\footnotetext{Słowa \code{for i = 1 to 10} oznaczają ,,dla $i=1$ do $10$'' --- przyp. tłum.}\nopagebreak

\bigpic{skrypt-kwadratowej-spirali}

Parametr bloku \code{przesuń} ma postać pomarańczowego owalu. Aby go tam umieścić, należy przeciągnąć taki sam owal będący częścią bloku \code{for}:\nopagebreak

\bigpic{przeciaganie-zmiennej}

Ten owal to \emph{zmienna} --- symboliczna nazwa reprezentująca jakąś wartość. Powyższy rysunek przedstawia sytuację sprzed zmiany drugiego parametru liczbowego bloku \code{for} z~domyślnego $10$ na $200$ oraz przeciągnięcia do jego środka bloku \code{obróć}. Blok \code{for} wykonuje swój parametr skryptowy wielokrotnie, podobnie jak \code{powtarzaj}, lecz przed każdym razem zapisuje liczbę do zmiennej~\code{i}, zaczynając od swojego pierwszego parametru liczbowego, dodając~$1$ przy każdym powtórzeniu, aż dojdzie do liczby z~drugiego parametru liczbowego. W~tym przypadku będziemy mieć $200$ powtórzeń, najpierw dla $\code{i}=1$, potem dla $\code{i}=2$, następnie $3$ i~tak dalej, aż do $\code{i}=200$ w~ostatnim powtórzeniu. W~rezultacie każdy blok \code{przesuń} rysuje coraz to dłuższy segment łamanej, co nadaje jej wygląd zbliżony do spirali. (Możesz spróbować ze skrętem $90$~stopni zamiast $92$; zobaczysz wtedy, dlaczego nazywamy tego rodzaju obraz ,,kwadratową spiralą'').

Zmienna~\code{i}~została utworzona przez blok \code{for} i~może zostać użyta wyłącznie wewnątrz jego klamry. Nawiasem mówiąc, jeśli nie spodoba nam się nazwa~\code{i}, możemy ją zmienić klikając pomarańczowy owal bez przeciągania go. Pokaże się wtedy okno dialogowe, do którego można wpisać inną nazwę:\nopagebreak

\bigpic{dialog-nazwy-zmiennej-skryptu}

Nazwa~,,\code{i}'' nie mówi nam zbyt wiele; można by tu użyć słowa ,,długość'', aby podkreślić znaczenie zmiennej. Nazwa~,,\code{i}'' jest popularna, gdyż w~matematyce istnieje tradycja używania liter od~\code{i} do~\code{n} dla liczb całkowitych. W~językach programowania nie musimy się jednak ograniczać do jednoliterowych nazw zmiennych.

\subsection{Zmienne globalne}

Możemy ,,ręcznie'' tworzyć zmienne, których widoczność nie jest ograniczona do jednego bloku. Aby to zrobić, należy kliknąć przycisk ,,\code{Stwórz zmienną}'' w~górnej części palety ,,Dane'':\nopagebreak

\bigpic{stworz-zmienna}

Otworzy się okno dialogowe, które pozwala nadać nazwę nowej zmiennej:\nopagebreak

\bigpic{dialog-nazwy-zmiennej}

Okno to pozwala również wybrać, czy zmienna ma być dostępna dla wszystkich duszków (co jest pożądane w~większości przypadków), czy ma być ona widoczna tylko dla aktualnego duszka. Postępujemy tak, jeśli mamy zamiar dać wielu duszkom własne zmienne o~\emph{tej samej nazwie}. Możemy potem kopiować skrypty między duszkami, przeciągając je z~obszaru skryptów aktualnego duszka na obrazek innego duszka w~zagrodzie duszków. Dzięki temu różne duszki będą wykonywać nieco inne rzeczy w~momencie wykonywania tego skryptu, ponieważ każdy z~nich będzie miał w~tej zmiennej zapisaną inną wartość.

Jeśli nadamy zmiennej nazwę ,,imię'', paleta ,,Dane'' będzie wyglądać następująco:\nopagebreak

\bigpic{zmienna-na-palecie}

Widzimy teraz przycisk ,,\code{usuń zmienną}'', a~także pomarańczowy owal z~nazwą zmiennej, taki sam jak owal na bloku \code{for}. Zmienną możemy przeciągnąć do dowolnego skryptu w~obszarze skryptów. Obok owalu widzimy pole wyboru; jest ono początkowo zaznaczone, dzięki czemu na scenie widoczny jest \emph{podgląd zmiennej}:\nopagebreak

\bigpic{podglad-zmiennej}

Kiedy nadamy zmiennej jakąś wartość, pojawia ona się w~polu podglądu.

Ale \emph{jak} nadać zmiennej wartość? Należy użyć do tego bloku \code{ustaw}:\nopagebreak

\bigpic{zapytaj-i-ustaw}

Zauważ, że \emph{nie przeciągamy} owalu zmiennej do bloku \code{ustaw}! Aby wybrać z~listy dostępnych zmiennych, należy kliknąć strzałkę w~dół przy pierwszym polu parametru tegoż bloku.

\subsection{Zmienne skryptu}

W~poprzednim przykładzie przeprowadzaliśmy interakcję z~użytkownikiem i~chcieliśmy zapamiętać jego imię na potrzeby całego projektu. To dobry przykład sytuacji, w~której właściwe jest użycie zmiennej \emph{globalnej} (z~rodzaju tych, które tworzymy przyciskiem ,,\code{Stwórz zmienną}''). Innym typowym przykładem jest zmienna ,,\code{wynik}'' w~projekcie gry. Czasem jednak potrzebujemy zmiennej tylko tymczasowo podczas wykonywania któregoś ze skryptów. W~takim przypadku możemy użyć bloku \code{zmienne skryptu} aby ją utworzyć:\nopagebreak

\begin{figure}[H]
\begin{minipage}{0.5\textwidth}
\centering
\includegraphics[scale=\defaultGraphicsScale]{skrypt-skrecajacej-linii}%
\end{minipage}%
\begin{minipage}{0.5\textwidth}
\centering
\includegraphics{../common/wiggling-line}
\end{minipage}%
\end{figure}

Podobnie jak w przypadku bloku \code{for}, aby zmienić nazwę zmiennej, wystarczy kliknąć pomarańczowy owal w bloku \code{zmienne skryptu} bez przeciągania. Możemy również stworzyć więcej tymczasowych zmiennych klikając strzałkę w prawo na końcu bloku. Spowoduje to dodanie kolejnego owalu zmiennej:\nopagebreak

\bigpic{zmienne-skryptu-a-b-c}

\section{I tak dalej}

Niniejszy podręcznik nie opisuje szczegółowo każdego bloku. Jest wiele bloków związanych z ruchem, dźwiękiem, kostiumami, efektami graficznymi i tak dalej. Ich przeznaczenie można poznać eksperymentalnie, ale także poprzez lekturę ,,ekranów pomocy''\footnote{Niestety, są one dostępne wyłącznie w języku angielskim --- przyp. tłum.}, które można obejrzeć klikając interesujący nas blok albo prawym przyciskiem myszy, albo lewym z wciśniętym jednocześnie klawiczem Ctrl, a następnie wybierając z menu polecenie ,,\code{pomoc\ldots}''. Jeśli zapomnisz, w której palecie znajduje się potrzebny Ci blok, ale pamiętasz choćby część nazwy, wciśnij Ctrl-F i wpisz ją w polu tekstowym, które pojawi się w obszarze palet.

\chapter{Zapisywanie i otwieranie projektów i multimediów}

Kiedy już stworzymy jakiś projekt, dobrze by było móc go zapisać.

After you've created a project, you'll want to save it, so that you can have access to it the next time you use \Snap{}. There are several ways to do that. You can save a project on your own computer, or you can save it at the \Snap{} web site. The advantage of saving on the net is that you have access to your project even if you are using a different computer, or a mobile device such as a tablet or smartphone. The advantage of saving on your computer is that you have access to the saved project while on an airplane or otherwise not on the net. This is why we have multiple ways to save.

\section{Local Storage}
\subsection{Localstore}
\subsection{XML Export}
\section{Cloud Storage}
\section{Loading Saved Projects}
\chapter{Building a Block}
\section{Simple Blocks}
\subsection{Custom Blocks with Inputs}
\section{Recursion}
\section{Block Libraries}
\chapter{First Class Lists}
\section{The list Block}
\section{Lists of Lists}
\section{Functional and Imperative List Programming}
\section{Higher Order List Operations and Rings}
\chapter{Typed Inputs}
\section{Scratch's Type Notation}
\section{The \Snap{} Input Type Dialog}
\subsection{Procedure Types}
\subsection{Pulldown inputs}
\subsection{Input variants}
\subsection{Prototype Hints}
\subsection{Title Text and Symbols}
\chapter{Procedures as Data}
\section{Call and Run}
\subsection{Call/Run with inputs}
\subsection{Variables in Ring Slots}
\section{Writing Higher Order Procedures}
\subsection{Recursive Calls to Multiple-Input Blocks}
\section{Formal Parameters}
\section{Procedures as Data}
\section{Special Forms}
\subsection{Special Forms in Scratch}
\chapter{Object Oriented Programming}
\section{Local State with Script Variables}
\section{Messages and Dispatch Procedures}
\section{Inheritance via Delegation}
\section{An Implementation of Prototyping OOP}
\chapter{The Outside World}
\section{The World Wide Web}
\section{Hardware Devices}
\section{Date and Time}
\chapter{Continuations}
\section{Continuation Passing Style}
\section{Call/Run w/Continuation}
\subsection{Nonlocal exit}
\chapter{User Interface Elements}
\section{Tool Bar Features}
\subsection{The \Snap{} Logo Menu}
\subsection{The File Menu}
\subsection{The Cloud Menu}
\subsection{The Settings Menu}
\subsection{Stage Resizing Buttons}
\subsection{Project Control Buttons}
\section{The Palette Area}
\subsection{Context Menus for Palette Blocks}
\subsection{Context Menu for the Palette Background}
\section{The Scripting Area}
\subsection{Sprite Appearance and Behavior Controls}
\subsection{Scripting Area Tabs}
\subsection{Scripts and Blocks Within Scripts}
\subsection{Scripting Area Background Context Menu}
\subsection{Controls in the Costumes Tab}
\subsection{The Paint Editor}
\subsection{Controls in the Sounds Tab}
\section{Controls on the Stage}
\section{The Sprite Corral and Sprite Creation Buttons}

\end{document}
